\chapter{Resistive Networks}
\section{Why They Matter}
Resistive networks set bias points, divide voltages, sample currents, terminate
transmission lines, and model losses. They are the simplest place to practice
KCL, KVL, and equivalent circuits.

\section{Series and Parallel}
Series resistors carry the same current:
\[
R_{\mathrm{eq}}=R_1+R_2+\cdots.
\]
Parallel resistors share the same voltage:
\[
\frac{1}{R_{\mathrm{eq}}}=\frac{1}{R_1}+\frac{1}{R_2}+\cdots.
\]
For two resistors,
\[
R_{\mathrm{eq}}=\frac{R_1R_2}{R_1+R_2}.
\]

\section{Voltage and Current Dividers}
For an unloaded voltage divider,
\[
V_o=V_i\frac{R_2}{R_1+R_2}.
\]
For two parallel resistors fed by total current \(I\),
\[
I_1=I\frac{R_2}{R_1+R_2},
\qquad
I_2=I\frac{R_1}{R_1+R_2}.
\]
Current divides inversely with resistance.

\section{Thevenin and Norton Equivalents}
Any linear one-port made of sources and resistors can be represented by a
Thevenin source \(V_{Th}\) in series with \(R_{Th}\), or a Norton source \(I_N\)
in parallel with \(R_N\):
\[
R_N=R_{Th},
\qquad
I_N=\frac{V_{Th}}{R_{Th}}.
\]
Thevenin and Norton models are especially useful when a load changes.

\begin{figure}[htbp]
\centering
\begin{circuitikz}
\draw (0,0) to[V,l=\(V_{Th}\)] (0,2)
      to[R,l=\(R_{Th}\)] (3,2)
      to[short,-o] (4,2);
\draw (0,0) -- (3,0) to[short,-o] (4,0);
\node at (4.8,1) {\(\Longleftrightarrow\)};
\draw (6,0) -- (6,2);
\draw (6,0) to[I,l_=\(I_N\)] (6,2);
\draw (6,2) to[R,l=\(R_N\)] (9,2) -- (9,0) -- (6,0);
\draw (9,2) to[short,-o] (10,2);
\draw (9,0) to[short,-o] (10,0);
\end{circuitikz}
\caption{Equivalent Thevenin and Norton one-port descriptions.}
\end{figure}

\section{Power Transfer and Matching}
For a purely resistive Thevenin source, maximum power transfer occurs when
\[
R_L=R_{Th}.
\]
For AC with complex impedances, maximum real power transfer uses the complex
conjugate match:
\[
Z_L=Z_S^*.
\]

\begin{exambox}
In RF systems, \(\SI{50}{\ohm}\) is a common reference impedance. Matching does
not mean all impedances in nature are \(\SI{50}{\ohm}\); it means the system has
been designed so sources, lines, loads, and measuring instruments interact in a
predictable way.
\end{exambox}

\section{Worked Example}
A \(\SI{12}{\volt}\) source feeds \(R_1=\SI{2.2}{\kilo\ohm}\) on top and
\(R_2=\SI{1.0}{\kilo\ohm}\) to ground. The unloaded output is
\[
V_o=12\frac{1.0}{2.2+1.0}=\SI{3.75}{\volt}.
\]
The Thevenin resistance seen by the load is
\[
R_{Th}=R_1\parallel R_2
=\frac{(2.2)(1.0)}{2.2+1.0}\,\si{\kilo\ohm}
=\SI{0.688}{\kilo\ohm}.
\]
